% Vlastni schematicke pudorysy; nejedna se o konkretni cizi PCB.
% Spolecny styl vlastnich vektorovych ilustraci; preklad jednotlivych stackup_*.tex.
\usepackage[czech]{babel}
\usepackage[utf8]{inputenc}
\usepackage[T1]{fontenc}
\usepackage{tgheros}
\renewcommand{\familydefault}{\sfdefault}
\definecolor{sig}{RGB}{244,185,75}
\definecolor{gnd}{RGB}{107,183,141}
\definecolor{pwr}{RGB}{228,129,112}
\definecolor{frfour}{RGB}{245,237,214}
\tikzset{copper/.style={draw=black!55,line width=0.4pt},
 layerlabel/.style={anchor=east,font=\small},
 note/.style={font=\small,align=center}}

\usetikzlibrary{arrows.meta}
\tikzset{signalpath/.style={red!75!black,line width=2.8pt,rounded corners=3pt},
 returnpath/.style={blue!75!black,line width=1.5pt,rounded corners=3pt,>=Latex},
 otherpath/.style={orange!85!black,line width=3pt}}
\newcommand{\SlotBoard}{
 \fill[gnd!65] (0,0) rectangle (9,6);
 \draw[black!60] (0,0) rectangle (9,6);
 \node[anchor=west,font=\small] at (0,6.4) {Pohled shora: F.Cu nad zelenou referenční GND na B.Cu};
}
\newcommand{\ConnectorExample}[1]{
 \begin{tikzpicture}[font=\small,>=Latex]
 \SlotBoard
 % Prekryv odstupu cizich PTH plosek vytvori souvisly kanal bez GND.
 \foreach \y in {.9,1.85,2.8,3.75,4.7} {
  \fill[white] (4,\y) circle (.53);
 }
 \foreach \y/\pin in {.9/1,1.85/2,2.8/3,3.75/4,4.7/5} {
  \filldraw[draw=black!60,fill=pwr] (4,\y) circle (.28);
  \fill[white] (4,\y) circle (.11);
  \node[anchor=east,font=\scriptsize] at (3.25,\y) {\pin};
 }
 \node[anchor=west,align=left,font=\scriptsize] at (5.7,.9) {Cizí PTH plošky konektoru\\nejsou připojené ke GND};
 \draw[->] (5.7,1.1)--(4.3,1.8);
 \ifnum#1=0
  \draw[signalpath] (.45,2.325)--(8.55,2.325);
  \draw[signalpath,->] (1,2.325)--(2,2.325);
  \draw[returnpath,->] (8.3,2.1)--(5.2,2.1)--(5.2,5.55)--(2.8,5.55)--(2.8,2.1)--(.7,2.1);
  \node[fill=white,inner sep=2pt,font=\scriptsize,align=center] at (7.1,4.5) {GND proud musí\\obejít řadu plošek};
  \node[anchor=east,align=right,font=\scriptsize] at (2.5,3.35) {Signál se mezi\\piny vejde};
  \draw[->] (2.7,3.3)--(3.9,2.325);
 \else
  \draw[signalpath] (.45,2.325)--(2.6,2.325)--(2.6,5.7)--(5.4,5.7)--(5.4,2.325)--(8.55,2.325);
  \draw[signalpath,->] (.8,2.325)--(1.8,2.325);
  \draw[returnpath,->] (8.3,2.1)--(5.6,2.1)--(5.6,5.4)--(2.35,5.4)--(2.35,2.1)--(.7,2.1);
  \node[fill=white,inner sep=2pt,font=\scriptsize,align=center] at (7.1,4.5) {Signál i návrat\\vedou stejnou oblastí};
 \fi
 \node[anchor=west,font=\scriptsize] at (0,-.45) {Červeně: signál na F.Cu. Modře: příklad vysokofrekvenčního návratu v GND.};
 \node[anchor=west,font=\scriptsize] at (0,-.8) {Bíle: oblast bez GND; nejde o vyfrézovaný otvor. Rozměry nejsou v měřítku.};
 \end{tikzpicture}
}
\newcommand{\CrossingExample}[1]{
 \begin{tikzpicture}[font=\small,>=Latex]
 \SlotBoard
 \fill[white] (3.58,.45) rectangle (4.42,5.2);
 \fill[white] (4,.45) circle (.42);
 \fill[white] (4,5.2) circle (.42);
 \draw[otherpath] (4,.45)--(4,5.2);
 \draw[signalpath] (.45,2.8)--(8.55,2.8);
 \draw[signalpath,->] (1,2.8)--(2,2.8);
 \node[anchor=west,align=left,font=\scriptsize] at (5.6,5) {Jiná síť na B.Cu\\+ její odstup od GND};
 \draw[->] (5.6,4.85)--(4.25,4.5);
 \ifnum#1=0
  \draw[returnpath,->] (8.3,2.6)--(4.95,2.6)--(4.95,5.8)--(3.05,5.8)--(3.05,2.6)--(.7,2.6);
  \node[fill=white,inner sep=2pt,font=\scriptsize,align=center] at (7.1,1.2) {Křížení není zkrat,\\ale GND je přerušená};
 \else
  % Tento GND mustek je na horni vrstve; nekoliduje se SIG na y=2.8.
  \draw[gnd!50!black,line width=4pt] (3.2,3.45)--(4.8,3.45);
  \foreach \x in {3.2,4.8} {
   \filldraw[fill=gnd!50!black,draw=black!60] (\x,3.45) circle (.16);
   \fill[white] (\x,3.45) circle (.07);
  }
  \draw[returnpath,->] (8.3,2.6)--(5.1,2.6)--(4.8,3.45)--(3.2,3.45)--(2.9,2.6)--(.7,2.6);
  \node[fill=white,inner sep=2pt,font=\scriptsize,align=center] at (6.8,1.1) {2 GND prokovy\\a krátký můstek na F.Cu};
  \draw[->] (5.7,1.6)--(4.8,3.25);
 \fi
 \node[anchor=west,font=\scriptsize] at (0,-.45) {Červeně: SIG na F.Cu. Oranžově: cizí trasa na B.Cu. Zeleně: GND na B.Cu.};
 \node[anchor=west,font=\scriptsize] at (0,-.8) {Modře: příklad návratu. Překryv tras v půdorysu neznamená jejich vodivé spojení.};
 \end{tikzpicture}
}
